\documentclass[10pt,a4paper]{amsart}
\usepackage{amsmath,amssymb,amsthm,graphicx}
\usepackage{verbatim}
\newenvironment{reference}{\comment}{\endcomment}
%\newenvironment{reference}{}{}
\newenvironment{slogan}{\comment}{\endcomment}
\newenvironment{history}{\comment}{\endcomment}

% For commutative diagrams we use Xy-pic
\usepackage[all]{xy}

% We use 2cell for 2-commutative diagrams.
\xyoption{2cell}
% Theorem environments.
%
\theoremstyle{plain}
\newtheorem{theorem}[subsection]{Theorem}
\newtheorem{proposition}[subsection]{Proposition}
\newtheorem{lemma}[subsection]{Lemma}

\theoremstyle{definition}
\newtheorem{definition}[subsection]{Definition}
\newtheorem{example}[subsection]{Example}
\newtheorem{exercise}[subsection]{Exercise}
\newtheorem{situation}[subsection]{Situation}

\theoremstyle{remark}
\newtheorem{remark}[subsection]{Remark}
\newtheorem{remarks}[subsection]{Remarks}

\numberwithin{equation}{subsection}

% Macros
%
\def\lim{\mathop{\mathrm{lim}}\nolimits}
\def\colim{\mathop{\mathrm{colim}}\nolimits}
\def\Spec{\mathop{\mathrm{Spec}}}
\def\Hom{\mathop{\mathrm{Hom}}\nolimits}
\def\Ext{\mathop{\mathrm{Ext}}\nolimits}
\def\SheafHom{\mathop{\mathcal{H}\!\mathit{om}}\nolimits}
\def\SheafExt{\mathop{\mathcal{E}\!\mathit{xt}}\nolimits}
\def\Sch{\mathit{Sch}}
\def\Mor{\mathop{\mathrm{Mor}}\nolimits}
\def\Ob{\mathop{\mathrm{Ob}}\nolimits}
\def\Sh{\mathop{\mathit{Sh}}\nolimits}
\def\NL{\mathop{N\!L}\nolimits}
\def\CH{\mathop{\mathrm{CH}}\nolimits}
\def\proetale{{pro\text{-}\acute{e}tale}}
\def\etale{{\acute{e}tale}}
\def\QCoh{\mathit{QCoh}}
\def\Ker{\mathop{\mathrm{Ker}}}
\def\Im{\mathop{\mathrm{Im}}}
\def\Coker{\mathop{\mathrm{Coker}}}
\def\Coim{\mathop{\mathrm{Coim}}}

% Boxtimes
%
\DeclareMathSymbol{\boxtimes}{\mathbin}{AMSa}{"02}

%
% Macros for moduli stacks/spaces
%
\def\QCohstack{\mathcal{QC}\!\mathit{oh}}
\def\Cohstack{\mathcal{C}\!\mathit{oh}}
\def\Spacesstack{\mathcal{S}\!\mathit{paces}}
\def\Quotfunctor{\mathrm{Quot}}
\def\Hilbfunctor{\mathrm{Hilb}}
\def\Curvesstack{\mathcal{C}\!\mathit{urves}}
\def\Polarizedstack{\mathcal{P}\!\mathit{olarized}}
\def\Complexesstack{\mathcal{C}\!\mathit{omplexes}}
% \Pic is the operator that assigns to X its picard group, usage \Pic(X)
% \Picardstack_{X/B} denotes the Picard stack of X over B
% \Picardfunctor_{X/B} denotes the Picard functor of X over B
\def\Pic{\mathop{\mathrm{Pic}}\nolimits}
\def\Picardstack{\mathcal{P}\!\mathit{ic}}
\def\Picardfunctor{\mathrm{Pic}}
\def\Deformationcategory{\mathcal{D}\!\mathit{ef}}
\usepackage{cleveref}
\crefname{originalsection}{original section}{original sections}
\Crefname{originalsection}{Original section}{Original sections}
\makeatletter
\newlabel{chow-section-degree-zero-cycles}{{0AZ0}{source}{}{sourcepointer-1}{}}
\newlabel{chow-section-degree-zero-cycles@cref}{{[originalsection][0][]0AZ0}{[1][0][]source}{}{}{}}
\makeatother
\crefname{originalreference}{original reference}{original references}
\Crefname{originalreference}{Original reference}{Original references}
\makeatletter
\newlabel{lemma-dim-0}{{0FHQ}{source}{}{sourcepointer-2}{}}
\newlabel{lemma-dim-0@cref}{{[originalreference][0][]0FHQ}{[1][0][]source}{}{}{}}
\makeatother
\crefname{originalreference}{original reference}{original references}
\Crefname{originalreference}{Original reference}{Original references}
\makeatletter
\newlabel{lemma-generated-by-separable}{{0FH5}{source}{}{sourcepointer-3}{}}
\newlabel{lemma-generated-by-separable@cref}{{[originalreference][0][]0FH5}{[1][0][]source}{}{}{}}
\makeatother
\title{The degree map on zero cycles agrees with the cohomological degree}
\author{The Stacks Project Authors}
\date{stacks-project:weil.tex:local-snapshot-20261012; GNU FDL 1.2 or later}
\begin{document}
\maketitle
\noindent\textit{Editorial note.} All delivered mathematical text is the author's own, copied from the pinned Stacks Project chapter file weil.tex; no statement or proof step was rewritten, reordered or removed.\par\medskip
\noindent\textit{Editorial note.} The counted proof is the author's complete argument: the verification on a closed point with separable residue field, and the reduction to generators of the group of zero cycles.\par\medskip
\noindent\textit{Editorial note.} The label lemma-degrees-cycles names the counted statement, and the proof environment immediately following it in the source is delivered verbatim as its proof.\par\medskip
\noindent\textit{Editorial note.} The counted statement is carried by the author's own lemma environment in the chapter source.\par\medskip
\noindent\textit{Editorial note.} Pre-open gate: python operations/agent-authoring/tools/precard\_gate.py --source weil.tex --label lemma-degrees-cycles --start 3504 --end 3535 --local-id 28157 returned RESULT=OK (L28157 free, occupancy FREE, public interval clear). Pre-assembly tag self-check: canonical tag resolved and not present in reference/public-items.jsonl.\par\medskip
\noindent\textit{Editorial note.} The statement and proof are copied from the pinned chapter file; the cited inputs are kept as external pointers with their canonical Stacks tags.\par\medskip
\noindent\textit{Allowed standard prerequisites (editorial).}\begin{itemize}
\item The degree of a zero dimensional scheme is computed by the trace. Same chapter, the computation of the degree map of a zero dimensional smooth projective scheme as the trace of the corresponding separable algebra. Identifies the cohomological degree of the class of a point with the degree of the residue field extension.
\item Generation of the group of zero cycles by classes of points with separable residue field. Same chapter, the lemma that the group of zero cycles is generated by classes of closed points with separable residue field. Reduces the commutativity of the diagram to the case of a single such point.
\item Compatibility of the cycle class map with pushforwards along a closed immersion. Axiom (C)(b) on the compatibility of the cycle class map with proper pushforwards. Transfers the class of the point to the cohomology of the scheme.
\end{itemize}\medskip
\ifdefined\hypertarget\hypertarget{sourcepointer-1}{}\fi
\noindent\textit{Source reference.} Original section 0AZ0: sources/stacks-chow/tags.txt, source lines 13808--13808. Prerequisite pointer invoked at proof line(s) 3518; it names a standard fact of a sibling chapter that the counted argument uses, and every step of the counted proof built on it is retained verbatim in the delivered body.\par\medskip
\ifdefined\hypertarget\hypertarget{sourcepointer-2}{}\fi
\noindent\textit{Source reference.} Original reference 0FHQ: sources/stacks-weil/pinned-original-source.tex, source lines 3432--3432. Prerequisite pointer invoked at proof line(s) 3532; it names a standard fact of the same chapter that the counted argument uses, and every step of the counted proof built on it is retained verbatim in the delivered body.\par\medskip
\ifdefined\hypertarget\hypertarget{sourcepointer-3}{}\fi
\noindent\textit{Source reference.} Original reference 0FH5: sources/stacks-weil/pinned-original-source.tex, source lines 2233--2233. Prerequisite pointer invoked at proof line(s) 3534; it names a standard fact of the same chapter that the counted argument uses, and every step of the counted proof built on it is retained verbatim in the delivered body.\par\medskip
\begin{lemma}
\label{lemma-degrees-cycles}
Assume given (D0), (D1), (D2), and (D3) satisfying (A), (B), and (C).
Let $X$ be a nonempty smooth projective scheme
equidimensional of dimension $d$ over $k$. The diagram
$$
\xymatrix{
\CH^d(X) \ar[r]_-\gamma \ar@{=}[d] &
H^{2d}(X)(d) \ar[d]^{\int_X} \\
\CH_0(X) \ar[r]^\deg & F
}
$$
commutes where $\deg : \CH_0(X) \to \mathbf{Z}$ is the degree of
zero cycles discussed in Chow Homology, Section
\ref{chow-section-degree-zero-cycles}.
\end{lemma}

\begin{proof}
Let $x$ be a closed point of $X$ whose residue field is separable
over $k$. View $x$ as a scheme and denote $i : x \to X$ the inclusion morphism.
To avoid confusion denote $\gamma' : \CH_0(x) \to H^0(x)$ the cycle class map
for $x$. Then we have
$$
\int_X \gamma([x]) = \int_X \gamma(i_*[x]) =
\int_X i_*\gamma'([x]) = \int_x \gamma'([x]) = \deg(x \to \Spec(k))
$$
The second equality is axiom (C)(b) and the third equality is
the definition of $i_*$ on cohomology. The final equality is
Lemma \ref{lemma-dim-0}. This proves the lemma
because $\CH_0(X)$ is generated by the classes of points $x$ as above
by Lemma \ref{lemma-generated-by-separable}.
\end{proof}

\end{document}
